\chapter{Debt Must Be Carried by an Asset}
\chaptermark{Debt Must Be Carried by an Asset}

Barbara Corcoran recalls buying a twelve-unit country motel whose papers
presented an attractive rent roll. After taking ownership, she discovered that
the scheduled rent had not been collected for two years. She later demolished
the property and absorbed the loss.

The document described what occupants were supposed to pay. The debt, taxes,
insurance, repairs, and ownership costs required cash that actually arrived.
That difference can ruin a leveraged purchase before any long-run thesis has
time to become true.

The title of this chapter is therefore shorthand, not magic. A building does
not make a payment. A truck, company, patent, or piece of equipment does not
make one either. Debt is carried by collected cash from a functioning asset,
supported by reserves and a borrower able to respond when the cash arrives
late. The asset matters because it may produce that cash and provide collateral.
Neither quality makes the payment automatic.

\section{Begin with the cash, not the price}

Ben Mallah describes a familiar property structure: the buyer contributes part
of the price and a bank lends the rest after examining the hotel and its
economics. In his rough ten-million-dollar example, the buyer supplies two
million and the bank eight million. He says the bank takes the risk.

The bank does take credit and collateral risk. The borrower may still lose the
entire equity contribution, further capital, the property, transaction costs,
time, reputation, and anything pledged through a guarantee or cross-collateral
agreement. Workers, residents, suppliers, and communities can also bear losses
without having chosen the leverage. The capital stack allocates risk; it does
not transfer all risk to the lender.

Start with cash generated by the asset under ordinary operation. For a property,
a simplified record is

\begin{equation}
\begin{gathered}
NOI=\text{collected operating revenue}\\
-\text{ordinary operating expense},\\[0.5ex]
CF_{\mathrm{owner}}=NOI-\text{debt service}\\
-\text{capital expenditure}\\
-\text{tax and other required cash}.
\end{gathered}
\end{equation}

The exact account definitions depend on the asset, agreement, and applicable
standards. Net operating income commonly excludes financing, income tax, and
capital expenditure, which is precisely why it should not be mistaken for cash
available to an owner. Use collected revenue, realistic vacancy or customer
loss, bad debt, concessions, management, maintenance, insurance, property tax,
utilities, compliance, and a reserve for irregular replacement. Do not
underwrite from a brochure's scheduled rent.

Then compare stressed operating cash with scheduled principal and interest:

\begin{equation}
DSCR_{\mathrm{stress}}
  =\frac{NOI_{\mathrm{stress}}}
         {\text{debt service}_{\mathrm{stress}}}.
\end{equation}

A ratio above one means the defined stressed NOI exceeds the defined debt
service. It does not fund a roof, lawsuit, tax shock, tenant relocation,
environmental repair, or family emergency omitted from the numerator. A
lender's minimum ratio is a covenant boundary, not a promise that the owner can
sleep.

\section{Write the complete capital stack}

A purchase price is only the first use of funds. Add diligence, legal work,
closing, immediate repair, equipment, working capital, initial vacancy,
interest before stabilization, and an operating reserve. Match those uses to
equity, senior debt, subordinate debt, a seller note, grants or incentives that
are actually available, and cash the operation can credibly produce.

\begin{equation}
\begin{gathered}
\text{purchase and required uses}\\
=\text{equity}+\text{debt}+\text{other committed sources}.
\end{gathered}
\end{equation}

For every layer of debt, record rate, amortization, payment dates, maturity,
balloon balance, collateral, lien priority, recourse, guarantees, covenants,
reporting duties, prepayment terms, default remedies, and whether the rate can
change. Record who supplied the equity, which rights it receives, and when more
capital can be required. A low opening payment can conceal a short maturity or
a large balloon.

This is the missing work behind the phrase \emph{other people's money}. Ken
McElroy describes banks and other institutions as channels through which
savers' capital reaches borrowers. Robert Kiyosaki celebrates the leverage and
also compares debt with a dangerous weapon. The second observation should
govern the first. Borrowed capital is another person's claim with priority,
terms, and remedies. Intelligence does not make that claim disappear; it makes
the borrower read it before signing.

\section{Let rent offset a home only after it is earned}

Mallah and Corcoran both suggest an owner-occupied small multiunit property:
live in one unit and collect rent from the others. The arrangement can reduce
the household's net housing cost while teaching the owner how to operate a
small asset. It can also combine home, investment, livelihood, debt, and local
concentration in one address.

The other units do not automatically pay all the bills. Begin with legal,
occupied, collected rent. Subtract vacancy, turnover, nonpayment, maintenance,
capital replacement, insurance, tax, utilities, management work, and financing.
Check zoning, habitability, fair-housing duties, leases, deposits, privacy,
safety, and the household's willingness to share the property with an
operating business.

The interviews name United States government-backed mortgage programs and
specific historical down payments. Eligibility, occupancy, property type,
insurance, limits, underwriting, fees, and rules change. A reader should verify
current terms with the official program and qualified local advisers rather
than carry an interview percentage into a purchase. A subsidy or guarantee can
alter financing; it cannot make a weak property affordable.

\section{Seller finance changes the lender, not the obligation}

Pace Morby describes buying from owners who prefer monthly payments to a cash
sale. The seller accepts a note and becomes the lender; rent from the property
is intended to fund the payments. In one reported 160-unit case, he gives
monthly gross income of about \$160,000 and about \$53,000 after the seller
payment, taxes, insurance, management, and repairs.

Seller finance can solve a real coordination problem. A seller may value income,
timing, price, or transition differently from a bank. A buyer may understand an
asset that conventional underwriting rejects. The parties can negotiate rate,
amortization, maturity, collateral, cure rights, and operating transition around
the actual property.

Removing a bank removes neither underwriting nor default. Verify title, liens,
authority, property condition, legal use, residents and leases, collections,
operating statements, tax, insurance, environmental and safety duties, deferred
maintenance, and the seller's right to foreclose. Model vacancy, repair, rate or
payment changes, a balloon, and the inability to refinance. Obtain independent
legal, tax, accounting, and technical review appropriate to the jurisdiction.
The seller's preferred tax timing is not the buyer's evidence of value.

Morby's claim that real estate is guaranteed should be rejected. Property can
lose tenants, access, insurability, lawful use, physical condition, market
value, or financing. A structure with no buyer cash at closing can expose the
buyer to guarantees, deficiency claims, unpaid work, legal duties, and losses
far beyond the initial cash contribution. Low cash in is a financing fact, not
a measure of safety.

\section[Refinancing and debt]{Refinancing returns capital by adding a claim}

Todd Nepola reports buying his first property for \$575,000, re-tenanting and
improving it, then refinancing against a reported value of \$800,000. He says a
\$600,000 loan returned his invested cash and closing costs while the property
continued to produce about \$5,000 a month. The full original debt, repair cost,
fees, tax, reserves, and cash-flow calculation are absent, so his return cannot
be reconstructed from the interview.

The mechanism can still be stated cleanly:

\begin{equation}
\begin{gathered}
\text{cash released}\\
=\text{new loan}-\text{old debt repaid}\\
-\text{closing cost}-\text{required reserve},\\[0.5ex]
\text{equity after refinance}\\
=\text{supported asset value}-\text{new debt}.
\end{gathered}
\end{equation}

Cash released is not profit created by the refinance. It is cash borrowed
against value that already exists or is believed to exist. Re-tenanting,
repair, better operation, market movement, or a more generous appraisal may
have raised the supported value; the new loan increases the fixed claim and may
return equity for use elsewhere.

Refinancing is productive when the remaining asset can carry the new terms
through a serious stress case and the released capital is used deliberately.
It becomes extraction when the owner removes the cushion that tenants,
customers, maintenance, and lenders rely on. Tax treatment is transaction- and
jurisdiction-specific and does not change the repayment obligation.

\section{The maturity date arrives even when the thesis is right}

Justin Waller describes underwriting net operating income, capitalization rate,
debt service, and cycle position, then imagines an adjustable loan moving from
roughly eight percent toward sixteen percent at the end of five to seven years.
Sixteen percent is his stress case, not a forecast. The useful question is why
the decision cannot wait until the reset.

Long before maturity, calculate the balance that must be repaid. Revalue the
asset from stressed, supportable income rather than hoped-for appreciation.
Apply a lower refinancing loan-to-value ratio, a higher interest rate, required
capital work, and realistic transaction cost. If the replacement loan is
smaller than the maturing balance, name the source of the gap. Sale is not a
complete answer unless price, time, buyer financing, tax, and closing risk are
also stressed.

The same test applies to an operating-company acquisition. A golf-course
interviewee reports buying a landscape-supply company, adding trucks, people,
and equipment, and later selling it. He recommends a government-backed
acquisition loan with a small down payment. The reported program amounts may be
dated or incomplete and should not be reused as present terms. A government
guarantee may protect part of a lender's loss; it does not generally erase the
borrower's repayment duty, collateral, guarantees, eligibility, fees, or
underwriting.

Additional equipment should carry its own financing through incremental
contribution after utilization, labor, fuel, maintenance, insurance, downtime,
and debt service. Buying the truck because the company is growing reverses the
logic. The evidence must show that the work can use the truck and collect before
the payment is due.

\section{Run the leverage survival test}
\enlargethispage{2\baselineskip}

Choose one property, company, equipment purchase, or other productive asset.
Test the financing as if neither appreciation nor easy refinancing arrives to
rescue it.

\begin{enumerate}[itemsep=0.14em]
\item Verify the right being bought, legal use, physical condition, operating
  history, counterparties, and every duty that travels with ownership.
\item Reconcile scheduled revenue to occupied, delivered, invoiced, and
  collected cash. Investigate every gap.
\item Build ordinary and stressed operating statements with vacancy or demand
  loss, bad debt, direct cost, management, maintenance, insurance, tax,
  compliance, and capital replacement.
\item Write all uses and committed sources. Record every debt layer, maturity,
  rate reset, balloon, lien, covenant, guarantee, recourse, and default remedy.
\item Calculate stressed debt service, coverage, cash after required spending,
  and the lowest cash balance by date.
\item Model a simultaneous shock: lower income, an expensive repair, delayed
  collection, higher rate, lower value, and tighter refinancing.
\item Name the cash, sale, equity, operating change, or negotiated repair that
  closes each gap. Reject a plan whose only repair is another uncommitted loan.
\item Set household and stakeholder boundaries: no hidden payroll delay,
  unsafe deferral, resident neglect, customer abandonment, or family guarantee
  beyond a loss consciously accepted in advance.
\item Decide the conditions to buy, renegotiate, reduce leverage, add reserves,
  bring in equity, stage the investment, or walk away.
\end{enumerate}

Debt can accelerate a useful asset and preserve capital for other work. It can
also turn one forecasting error into a forced sale at the worst possible time.
The disciplined borrower does not ask whether debt is good. The borrower asks
whether the cash, structure, reserve, and people exposed to failure can carry
this debt through the life it is actually likely to have.

One asset, however well financed, remains one uncertain outcome. The next task
is to let time and diversification build resilience without dissolving every
advantage into an average.
